\documentclass[a4paper,12pt]{article}
\usepackage[utf8]{inputenc}
\usepackage[T1]{fontenc}
\usepackage[french]{babel}
\usepackage{amsmath,amssymb}
\usepackage{geometry}
\usepackage{tikz}

\geometry{margin=2.5cm}
\newcommand{\degmark}{{}^{\circ}}
\title{Probl\`eme de G\'eom\'etrie -- Trouver l'angle $\alpha$}
\author{S.\ Jaubert}
\date{F\'evrier 2026}

\begin{document}
\maketitle

\section*{\'Enonc\'e}

Soit $ABC$ un triangle avec $\angle ABC = 3\alpha$ et $\angle BAC = \alpha$.
Le point $M$ est le milieu de $[AB]$ ($AM = BM$).
La m\'ediane $[CM]$ forme un angle $\angle CMA = 4\alpha$.

Trouver $\alpha$.

\begin{center}
\begin{tikzpicture}[scale=3.5]
  % alpha = 32.594 deg
  \pgfmathsetmacro{\adeg}{32.594}
  \pgfmathsetmacro{\Bdeg}{3*\adeg}  % angle at B = 97.78
  \pgfmathsetmacro{\Adeg}{\adeg}     % angle at A = 32.59
  \pgfmathsetmacro{\Cdeg}{180 - 4*\adeg} % angle at C = 49.62

  % Place B at origin, A at (1,0)
  \coordinate (B) at (0,0);
  \coordinate (A) at (1,0);
  \coordinate (M) at (0.5,0);

  % C : from B, direction 3*alpha from BA
  % BC = sin(alpha)/sin(4*alpha)
  \pgfmathsetmacro{\BC}{sin(\adeg)/sin(4*\adeg)}
  \coordinate (C) at ({\BC*cos(\Bdeg)},{\BC*sin(\Bdeg)});

  % Draw triangle
  \draw[thick] (B)--(A)--(C)--cycle;
  \draw[thick] (C)--(M);

  % Extend base line
  \draw[thin,gray] (-0.15,0)--(1.15,0);

  % Labels
  \node[below left] at (B) {$B$};
  \node[below right] at (A) {$A$};
  \node[above] at (C) {$C$};
  \node[below] at (M) {$M$};

  % Tick marks for AM = BM
  \draw (0.22,0.03)--(0.28,-0.03);
  \draw (0.72,0.03)--(0.78,-0.03);

  % Angle at B = 3*alpha
  \draw (0.12,0) arc[start angle=0, end angle=\Bdeg, radius=0.12];
  \node at ({0.2*cos(\Bdeg/2)},{0.2*sin(\Bdeg/2)}) {$3\alpha$};

  % Angle at A = alpha
  \pgfmathsetmacro{\startA}{180 - \adeg}
  \draw ({1+0.15*cos(\startA)},{0.15*sin(\startA)}) arc[start angle=\startA, end angle=180, radius=0.15];
  \node at ({1+0.22*cos(\startA/2 + 90)},{0.22*sin(\startA/2 + 90)}) {$\alpha$};

  % Angle CMA = 4*alpha at M
  \pgfmathsetmacro{\angleCM}{atan2(\BC*sin(\Bdeg)-0, \BC*cos(\Bdeg)-0.5)}
  \draw ({0.5+0.1*cos(0)},{0.1*sin(0)}) arc[start angle=0, end angle=\angleCM, radius=0.1];
  \node at ({0.5+0.16*cos(\angleCM/2)},{0.16*sin(\angleCM/2)}) {$4\alpha$};

  % AM = BM annotation
  \node at (0.85,0.35) {$AM = BM$};
\end{tikzpicture}
\end{center}

\section*{Solution}

\textbf{\'Etape 1 -- Angles des sous-triangles.}

Puisque $M \in (AB)$, les angles $\angle CMA$ et $\angle CMB$ sont suppl\'ementaires~:
\[
  \angle CMB = \pi - 4\alpha.
\]

Dans le \textbf{triangle $BCM$}~:
\[
  \angle MBC + \angle BMC + \angle BCM = \pi
  \quad\Longrightarrow\quad
  3\alpha + (\pi - 4\alpha) + \angle BCM = \pi
  \quad\Longrightarrow\quad
  \angle BCM = \alpha.
\]

Dans le \textbf{triangle $ACM$}~:
\[
  \angle MAC + \angle CMA + \angle ACM = \pi
  \quad\Longrightarrow\quad
  \alpha + 4\alpha + \angle ACM = \pi
  \quad\Longrightarrow\quad
  \angle ACM = \pi - 5\alpha.
\]

\textbf{Condition de validit\'e :} $\pi - 5\alpha > 0$, soit $\alpha < 36\degmark$.

\bigskip

\textbf{\'Etape 2 -- Loi des sinus dans chaque triangle.}

On pose $a = BM = MA$.

\medskip
\emph{Triangle $BCM$} (angles : $3\alpha$ en $B$,\; $\pi - 4\alpha$ en $M$,\; $\alpha$ en $C$)~:
\[
  \frac{a}{\sin(\alpha)} = \frac{CM}{\sin(3\alpha)}
  \qquad\Longrightarrow\qquad
  CM = \frac{a\,\sin(3\alpha)}{\sin(\alpha)}.  \tag{1}
\]

\emph{Triangle $ACM$} (angles : $\alpha$ en $A$,\; $4\alpha$ en $M$,\; $\pi - 5\alpha$ en $C$)~:
\[
  \frac{a}{\sin(\pi - 5\alpha)} = \frac{CM}{\sin(\alpha)}
  \qquad\Longrightarrow\qquad
  CM = \frac{a\,\sin(\alpha)}{\sin(5\alpha)}.  \tag{2}
\]

\bigskip

\textbf{\'Etape 3 -- \'Equation trigonom\'etrique.}

En \'egalant (1) et (2)~:
\[
  \frac{\sin(3\alpha)}{\sin(\alpha)} = \frac{\sin(\alpha)}{\sin(5\alpha)}
  \qquad\Longrightarrow\qquad
  \boxed{\sin(3\alpha)\,\sin(5\alpha) = \sin^2(\alpha)}.
\]

\bigskip

\textbf{\'Etape 4 -- Lin\'earisation (formules de produit).}

\begin{itemize}
  \item \`A gauche~: $\displaystyle \sin(3\alpha)\sin(5\alpha)
    = \tfrac{1}{2}\bigl[\cos(2\alpha) - \cos(8\alpha)\bigr]$.
  \item \`A droite~: $\displaystyle \sin^2(\alpha)
    = \tfrac{1}{2}\bigl[1 - \cos(2\alpha)\bigr]$.
\end{itemize}

L'\'equation devient~:
\[
  \cos(2\alpha) - \cos(8\alpha) = 1 - \cos(2\alpha)
\]
\[
  \boxed{\cos(8\alpha) = 2\cos(2\alpha) - 1}.
\]

\bigskip

\textbf{\'Etape 5 -- R\'eduction polynomiale.}

On pose $t = \cos(2\alpha)$. Le polyn\^ome de Tchebychev donne
$\cos(8\alpha) = T_4(t) = 8t^4 - 8t^2 + 1$. Donc~:
\[
  8t^4 - 8t^2 + 1 = 2t - 1
  \qquad\Longrightarrow\qquad
  4t^4 - 4t^2 - t + 1 = 0.
\]

$t = 1$ est racine \'evidente (correspondant \`a $\alpha = 0$, trivial). On factorise~:
\[
  (t - 1)\underbrace{(4t^3 + 4t^2 - 1)}_{Q(t)} = 0.
\]

\emph{V\'erification~:} $(t-1)(4t^3+4t^2-1) = 4t^4 + 4t^3 - t - 4t^3 - 4t^2 + 1 = 4t^4 - 4t^2 - t + 1$.\;\checkmark

\bigskip

\textbf{\'Etape 6 -- R\'esolution du cubique.}

Le polyn\^ome $Q(t) = 4t^3 + 4t^2 - 1$ n'a qu'une seule racine r\'eelle dans $[-1,\;1]$.

Son discriminant vaut $\Delta = -4(4)^3(-1) - 27(4)^2(-1)^2 = 256 - 432 = -176 < 0$,
ce qui confirme une seule racine r\'eelle (les deux autres sont complexes conjugu\'ees).

Par la m\'ethode de Cardan ou par r\'esolution num\'erique~:
\[
  t_0 = \cos(2\alpha) \approx 0{,}4196
  \qquad\Longrightarrow\qquad
  2\alpha \approx 65{,}19\degmark
\]
\[
  \boxed{\alpha \approx 32{,}59\degmark}
\]

\bigskip

\textbf{\'Etape 7 -- V\'erification.}

\begin{center}
\begin{tabular}{ll}
\hline
Angle & Valeur \\
\hline
$\angle BAC = \alpha$ & $32{,}59\degmark$ \\
$\angle ABC = 3\alpha$ & $97{,}78\degmark$ \\
$\angle BCA = \pi - 4\alpha$ & $49{,}62\degmark$ \\
$\angle CMA = 4\alpha$ & $130{,}38\degmark$ \\
$\angle ACM = \pi - 5\alpha$ & $17{,}03\degmark$ \\
$\angle BCM = \alpha$ & $32{,}59\degmark$ \\
\hline
Somme dans $\triangle ABC$ & $180\degmark$ \;\checkmark \\
\hline
\end{tabular}
\end{center}

V\'erification num\'erique de l'\'equation trigonom\'etrique~:
\[
  \sin(3\times 32{,}59\degmark)\,\sin(5\times 32{,}59\degmark)
  \approx 0{,}2899
  = \sin^2(32{,}59\degmark) \approx 0{,}2899.
  \quad\checkmark
\]

\bigskip

\textbf{Remarque.}
La valeur exacte de $\alpha$ est donn\'ee par
\[
  \alpha = \frac{1}{2}\arccos(t_0)
\]
o\`u $t_0$ est l'unique racine r\'eelle du cubique $4t^3 + 4t^2 - 1 = 0$.
Cette valeur n'est pas un multiple rationnel de $\pi$, ni un angle
remarquable en degr\'es.

\end{document}
